\documentclass[instructions.tex]{subfiles}
\begin{document}
 
\chapter{De la rigueur des peines de l'Enfer}
 
L'enfer est un si grand malheur, que de tous les hommes, quand il n'y en aurait qu'un seul qui dût y être condamné, nous devrions tous craindre d'y tomber. Deux sortes de peines sont en enfer : celle du sens, principalement causée par le feu; et la peine du dam, causée par la perte de Dieu. L'une et l'autre sont incompréhensibles.

\primo Il n'y a point de proportion entre le feu de l'enfer et celui des fournaises embrasées; les feux d'ici-bas, en comparaison, ne sont que comme une ombre; il fait souffrir tous les tourmens; il fait souffrir également l'ame dans toutes ses puissances, et le corps dans tous ses sens. Il n'y a point de supplice qu'un réprouvé n'éprouve\footnote{Omnis dolor simul super eum. Job. 20. 22.}. Les Néron, les Dioclétien, tout ce qu'il y a jamais eu de monstres en cruauté, n'ont rien inventé qui soit comparable. C'étaient des hommes, et ici c'est un Dieu qui punit en Dieu, et qui dans sa fureur emploie sa toute-puissance\footnote{Ignis succensus est in furore meo. Jac.15.14.}.

L'esprit humain ne peut rien imaginer de plus horrible que l'état d'un réprouvé. Représentez-vous des corps embrasés, infects, couverts de lèpre, de cancers, de pourriture, plongés dans un gouffre de feu, de soufre brûlant; des corps vivans déchirés, écorchés, disloqués, entassés les uns sur les autres, dont l'odorat, la bouche, les yeux, tous les membres, tous les sens, souffrent en particulier le plus violent supplice, sans adoucissement et sans relâche\footnote{Congregabo super eso mala 32.32}. 

Leurs cris sont si perçans, qu'ils seraient capables de fendre les rochers; leurs douleurs, si cuisantes et si vives, qu'une heure de ces tourmens est plus insupportable que vingt ans en ce monde de la maladie la plus aiguë; la puanteur de ces corps est si horrible, que saint Bonaventure assure qu'un seul serait capable d'infecter l'univers.

Tout ce qui se présente aux les tourmente. Ils ne touchent que du feu; ils ne sentent que du feu; ils ne voient que des objets pénétrés de feu. Tout les remplit d'horreur et d'effroi; ils ne voient que des figures hideuses, des monstres horribles, des démons cruels sous des formes épouvantables; ils n'entendent que des lamentations, d'horribles cris de rage, des grincemens de dents\footnote{Ibi erit fletus et stridor dentium. Matth. 22.13.}.

Voilà le tombeau des délices du monde, le terme de la volupté et de l'ambition. Oh! que la délicatesse avec laquelle on flate aujourd'hui son corps, sera cruellement payée un jour! que ses plaisirs, ses intempérances, lui coûteront cher!

Vous frémissez à cette seule peinture de l'enfer; hélas! ce n'est qu'une légère ébauche de ce qu'on y souffre. Si la gloire du ciel est au-dessus de nos pensées, les tourmens de l'enfer le sont aussi; et nous montre l'intelligence; l'esprit de l'homme est trop borné pour se former une idée de la justice du Tout-Puissant. Lorsqu'il châtia l'Egypte par des fléaux qu'il a désolée, ce n'était, dit l'Ecriture, que le doigt de Dieu qui frappait; mais, pour punir Lucifer et les réprouvés, il emploie sa toute-puissance et la force de son bras\footnote{Fecit potentiam in brachio suo. Luc. 1.51.}.

Malheur à vous, si vous attendez de le croire lorsqu vous l'éprouverez.

Si vous craignez les disgraces et les maux de cette vie, qui durent peu de temps, pourquoi ne prévenez-vous pas des maux extrêmes qui ne finiront jamais? vous auriez peine de souffrir la piqûre d'une abeille, comment souffrirez-vous l'activité d'un feu dévorant? Que votre aveuglement est déplorable si vous ne pensez pas au terme où doit venir votre attachement aux plaisirs, aux vanités, aux biens de la terre, et l'abus que vous faites des graces du ciel!

L'homme sage agit bien autrement; il sait qu'on ne peut être heureux en ce monde et en l'autre; jouir de tous les plaisirs de la terre et de ceux du ciel; c'est tous les jours qu'il ne s'attache point aux biens et aux plaisirs de ce monde; qu'il tâche de s'assurer ceux du ciel; qu'il pleure ses fautes, qu'il fait pénitence, pour ne pas pleurer éternellement. Je châtie mon corps, disait saint Paul, pour ne pas devenir un réprouvé. Il ne suis réduit comme un prisonnier dans un désert; disait saint Jérôme, de peur de tomber dans l'enfer, et de me perdre avec les mondains\footnote{Ob gehennae metum, tali mecarcere damnaveram. S.Hier.}.

C'est donc avec raison que Gui le Chartreux a dit que, si la pénitence et les tribulations de cette vie sont la semence d'une joie et d'une gloire immortelle; les plaisirs du temps, l'attache aux biens, sont aussi une semence de douleurs et de regrets, qui poussent leur germe dans l'éternité\footnote{Tria temporalia gaudia futuraorum sunt semina dolorum. Train. c. 4.}.

Regrets désespérans, qui, comme un ver rongeur, déchirent sans relâche le cœur d'un réprouvé\footnote{Vermis eorum non moritur. Marc. 9.}! Je pouvais gagner le ciel; j'en ai eu le temps et les moyens, mais j'en ai abusé : de quoi me servent à présent les plaisirs, tout ce que j'ai fait et possédé sur la terre\footnote{Quid profuit nobis?}?

\secundo Mais la peine la plus vive et la plus cruelle d'une ame damnée, c'est d'être séparée de Dieu, de savoir que jamais elle ne jouira de son Dieu. J'étais faite pour Dieu, et jamais je ne le verrai; il devait faire mon bonheur, et jamais je ne le posséderai; je l'ai perdu par ma faute; je l'ai perdu pour peu de chose, je l'ai perdu sans ressource. J'AI PERDU DIEU! O pensée accablante! et que le supplice qu'elle cause est insupportable!

Nous ne comprenons pas en cette vie ce que c'est que Dieu, et ce que c'est que de le perdre; mais le réprouvé le comprend, et ressent si cruellement la perte qu'il en a faite, que si, au milieu de ses brasiers, il pouvait espérer de voir son Dieu, et de le posséder pendant une heure après chaque mille ans, il serait content et souffrirait ses tourmens avec consolation.

A quoi pensez-vous, lorsque vous perdez Dieu en perdant sa grâce par le péché mortel? Et vous, ames tièdes, que faites-vous, lorsque, par tant de fautes légères, que vous traitez de bagatelles, vous vous exposez à tomber dans le péché mortel, et à perdre Dieu? les damnés pleurent la perte qu'ils ont faite de Dieu, et vous y êtes insensibles! Les damnés voudraient retourner à Dieu, mais il n'est plus temps, ils n le peuvent plus : vous en avez le temps, et vous le pouvez encore : vous ne le voulez pas; vous avez donc le cœur plus dur qu'un réprouvé. Cherchez Dieu dans le temps que vous pouvez le trouver; retournez à lui à présent qu'il vous rappelle; dans l'éternité vous ne le pourrez plus.

\section{Résolutions}

\primo Je gémirai de l'aveuglement que j'ai eu jusqu'à ce jour, de penser si peu aux châtimens épouvantables que j'ai mérité de souffrir en enfer. 

\secundo Je me dirai souvent : Si, par un miracle qui n'aura jamais lieu, un réprouvé était rappelé sur la terre, et que Dieu lui accordât le temps et les moyens que sa bonté veut encore me donner, pour réparer le passé et gagner le ciel, que ferait ce réprouvé? pour ce qu'il ferait, il faut que je le fasse moi-même, et dès ce jour. 

\prayer{O mon Dieu! préservez-moi de l'enfer.}
 
\end{document}
